\section{Related Work}

\paragraph{Conversational tutoring.}
CiTT's Socratic interaction design is related to intelligent tutoring systems and dialogue-based tutors such as AutoTutor, which demonstrated mixed-initiative natural-language tutoring as an educational interaction pattern \cite{graesser_autotutor_2001}. CiTT differs by using the dialogue as an interface to executable circuit models rather than as the final source of engineering authority.

\paragraph{LLMs for technical reasoning.}
LLMs can solve and explain many textbook problems, but recent work on mathematical reasoning emphasizes robustness and reliability limitations under perturbation, symbolic ambiguity, and multi-step reasoning \cite{llm_math_reasoning_survey}. CiTT does not treat language models as useless. Instead, it treats language reasoning as a front end that should be grounded in executable engineering artifacts for complex circuit tutoring.

\paragraph{Circuit simulation and MNA.}
Classical circuit simulation uses formulations such as modified nodal analysis and SPICE-style nonlinear circuit analysis \cite{ho_mna_1975,nagel_rohrer_1971}. CiTT's development initially explored a small internal MNA-style solver, but shifted to Simscape because general nonlinear biomedical circuits, visualization, and component coverage are better handled by mature simulation infrastructure.

\paragraph{Model-based design and physical modeling.}
Simulink provides a block-diagram environment for simulation and Model-Based Design, while Simscape supports physical modeling with component libraries and solver integration \cite{mathworks_simulink,mathworks_simscape,mathworks_simscape_electrical}. CiTT uses these tools as the engineering substrate for tutoring, rather than generating only prose or static diagrams.

\paragraph{Agentic MATLAB and Simulink workflows.}
The Simulink Agentic Toolkit gives AI agents tools and skills for working with Simulink models through MATLAB MCP Server and Model Context Protocol workflows \cite{mathworks_satk,matlab_agentic_toolkit}. CiTT builds on this direction by turning agentic model construction into a student-facing tutoring loop.

\paragraph{Biomedical instrumentation education.}
Biomedical instrumentation and voltage-clamp teaching motivate the benchmark design. Standard instrumentation texts emphasize sensor interfaces, amplifiers, filtering, and measurement systems \cite{webster_medical_instrumentation_2020}, while the voltage-clamp benchmark is grounded in the historical role of feedback experiments for measuring membrane current-voltage behavior \cite{hodgkin_huxley_katz_1952}. CiTT's contribution is educational workflow integration, not a new electrophysiology model.
